\chapter{Karamata's theorem}

Let $F \colon \bR \to \bR$ be antitone and regularly varying with index $-\alpha$, where
$\alpha > 1$. Typically $F = 1 - G$ is the survival function of a c.d.f.\ $G$ in the
Fréchet domain of attraction (Chapter~\ref{chap:domains-of-attraction}).

\begin{lemma}[Positivity of antitone eventually positive functions]
  \label{lem:pos-of-antitone}
  \lean{ExtremeValueTheory.IsRegularlyVarying.pos_of_antitone}
  \leanok
  An antitone function $F$ with $F(x) > 0$ for all sufficiently large $x$ is positive everywhere.
\end{lemma}
\begin{proof}
  \leanok
  For any $t$ there is $x \ge t$ with $F(x) > 0$, and $F(t) \ge F(x)$.
\end{proof}

\begin{lemma}[Integrability of regularly varying tails]
  \label{lem:rv-integrableOn-Ioi}
  \uses{def:regularly-varying, thm:rv-potter-monotone}
  \lean{ExtremeValueTheory.IsRegularlyVarying.integrableOn_Ioi}
  \leanok
  If $F$ is antitone and regularly varying with index $-\alpha$, $\alpha > 1$, then $F$ is
  integrable on every half-line $(x, \infty)$.
\end{lemma}
\begin{proof}
  \leanok
  By the Potter bounds (Theorem~\ref{thm:rv-potter-monotone}) with
  $\delta = (\alpha - 1)/2$, $F(t) \le C t^{-\alpha+\delta}$ for $t$ large, and
  $-\alpha + \delta < -1$, so $F$ is integrable near $+\infty$. On bounded intervals
  $F$ is bounded and measurable, being monotone.
\end{proof}

\begin{theorem}[Karamata's theorem for tail integrals]
  \label{thm:karamata}
  \uses{def:regularly-varying, thm:rv-potter-monotone, lem:rv-integrableOn-Ioi, lem:pos-of-antitone}
  \lean{ExtremeValueTheory.IsRegularlyVarying.tendsto_integral_Ioi_div}
  \leanok
  Let $F$ be antitone and regularly varying with index $-\alpha$, where $\alpha > 1$. Then
  \begin{align*}
    \lim_{x \to \infty} \frac{\int_x^\infty F(t) \, \ud t}{x \, F(x)} \; = \; \frac{1}{\alpha - 1}
  \end{align*}
  [BGT, Proposition~1.5.10; see also Theorem~1.5.11]. In particular, if $F$ is the survival
  function of a random variable, the mean excess over a high threshold $x$ is asymptotically
  $x/(\alpha - 1)$.
\end{theorem}
\begin{proof}
  \leanok
  After the substitution $t = x u$,
  \begin{align*}
    \frac{\int_x^\infty F(t) \, \ud t}{x \, F(x)} \; = \; \int_1^\infty \frac{F(x u)}{F(x)} \, \ud u .
  \end{align*}
  The integrand tends to $u^{-\alpha}$ pointwise. By the Potter bounds
  (Theorem~\ref{thm:rv-potter-monotone}) with $\delta = (\alpha-1)/2$ and $A = 2$, for $x$
  large it is dominated by $2 u^{-\alpha + \delta}$, which is integrable on $(1, \infty)$.
  Dominated convergence yields the limit $\int_1^\infty u^{-\alpha} \, \ud u = \frac{1}{\alpha - 1}$.
\end{proof}

\begin{theorem}[Karamata's theorem, relaxed hypotheses]
  \label{thm:karamata-relaxed}
  \uses{thm:karamata}
  \lean{ExtremeValueTheory.IsRegularlyVarying.tendsto_integral_Ioi_div_of_measurable,
    ExtremeValueTheory.IsRegularlyVarying.tendsto_integral_Ioi_div_of_antitoneOn}
  \leanok
  The conclusion of Theorem~\ref{thm:karamata} holds for every measurable $F \in RV_{-\alpha}$,
  $\alpha > 1$, and for every $F \in RV_{-\alpha}$ that is antitone on some $[a, \infty)$
  [BGT, Proposition~1.5.10].
\end{theorem}
\begin{proof}
  \leanok
  Potter bounds for measurable regularly varying functions give the domination needed in the
  proof of Theorem~\ref{thm:karamata}.
\end{proof}

\begin{theorem}[Karamata's theorem, converse half]
  \label{thm:karamata-converse}
  \uses{thm:karamata-representation}
  \lean{ExtremeValueTheory.IsRegularlyVarying.of_tendsto_mul_div_integral_Ioi, ExtremeValueTheory.IsRegularlyVarying.of_log_sub_eq_integral,
    ExtremeValueTheory.IsRegularlyVarying.of_tendsto_mul_deriv_div, ExtremeValueTheory.IsRegularlyVarying.of_abs_log_sub_le, ExtremeValueTheory.IsRegularlyVarying.of_tendsto}
  \leanok
  Let $F$ be antitone and integrable on $(a, \infty)$, and $\alpha > 1$. If
  $x F(x) / \int_x^\infty F \to \alpha - 1$, then $F \in RV_{-\alpha}$ [BGT, Theorem~1.6.1].
  Similarly, if $g > 0$ is differentiable with $x g'(x)/g(x) \to \rho$, then $g \in RV_\rho$.
\end{theorem}
\begin{proof}
  \leanok
  With $I(x) = \int_x^\infty F$, $(\log I)'(x) = -F(x)/I(x)$, so $x (\log I)'(x) \to 1 - \alpha$.
  The representation theorem gives $I \in RV_{1-\alpha}$, and $F = (I/x) \cdot (x F/I)$.
\end{proof}

\begin{theorem}[Karamata's theorem for integrals over $[0, x]$]
  \label{thm:karamata-zero}
  \uses{thm:karamata-relaxed}
  \lean{ExtremeValueTheory.IsRegularlyVarying.tendsto_integral_zero_div}
  \leanok
  If $f \in RV_\rho$ is measurable and locally integrable on $[0, \infty)$ and $\rho > -1$, then
  $\int_0^x f / (x f(x)) \to 1/(\rho + 1)$ [BGT, Proposition~1.5.8].
\end{theorem}
\begin{proof}
  \leanok
  Uniform convergence on compacts and Potter bounds near $0$ after rescaling $t = x s$.
\end{proof}
